\begin{abstract}
This thesis investigates whether constraints formulated within a Logic Tensor Network (LTN) perspective can improve hippocampus segmentation when labelled data and training duration are limited. The study concentrates on differentiable supervision around the annotated foreground boundary and examines its use in both a three-dimensional SwinUNETR model and parameter-efficient adaptation of MedSAM3. The main question concerns performance under a specified training policy; a separate question is whether an advantage persists with longer optimization.

In the audited SwinUNETR study, ten labelled case volumes per fold were used across four folds and three optimization seeds. Adding a calibrated boundary-band loss increased validation-selected foreground Dice by an average of 3.882 percentage points under a maximum-75-epoch policy. Exploratory longer-training controls on one fold reduced the observed difference. In a separate MedSAM3 pilot, fine-tuning on five labelled volumes for 200 optimizer updates increased mean whole-hippocampus validation Dice from 0.865080 to 0.877599 when boundary-band and edge-consistency losses were added. This pilot used one seed and five development-validation volumes; a broader replication is incomplete in this evidence snapshot.

These observations support a conditional benefit from boundary-focused supervision under restricted training schedules. They do not establish a general improvement in fully optimized performance or patient-independent generalization. The band objective is equivalent to a conventional balanced boundary binary cross-entropy loss, so the results concern the implemented inductive bias rather than the superiority of logical notation itself.
\end{abstract}
